%=============================================================================!
% 5.7  DEGREES OF FREEDOM                                                     !
%=============================================================================!
\section{Degrees of freedom}
\label{sec:ro-freedom}

The positions of $N$ atoms are $3N$ numbers, so a molecule can move in
$3N$ independent ways, called its degrees of freedom. Some of these move
the molecule as a whole, without changing any distance inside it. This
section counts them, shows that they appear as zero frequencies among the
normal modes of \cref{sec:os-modes}, and finds how many ways a molecule
has left to vibrate.

\subsection*{Counting}

A rigid motion of the whole molecule is a drift in some direction, a
combination of drifts along $x$, $y$ and $z$, three ways, plus a turning
about some axis through the centre of mass, a combination of turnings
about $x$, $y$ and $z$, three more. For a linear molecule, turning about
its own line moves no atom (\cref{sec:ro-remove}), so it has only two.
What is left changes the shape of the molecule:
\begin{equation*}
   3N - 6 \text{ internal degrees of freedom, or } 3N - 5
   \text{ for a linear molecule.}
\end{equation*}
Water, with $N = 3$, has $9 - 6 = 3$, the two stretches and the bend of
\cref{tab:os-vibrations}; carbon dioxide, linear, has $9 - 5 = 4$;
methane, CH$_4$, has $15 - 6 = 9$, and benzene, C$_6$H$_6$, $36 - 6 =
30$.

\subsection*{Rigid motions as zero modes}

Near its minimum the motion of a molecule is a sum of normal modes, with
squared angular frequencies the eigenvalues of the mass-weighted Hessian
(\cref{sec:os-modes}). A rigid motion stretches nothing, so we expect
nothing to pull it back: a mode of zero frequency. We show this for a
molecule held together by springs between pairs of atoms.

\emph{The Hessian of one spring.} A spring of stiffness $k$ joins atoms $i$
and $j$ and stores $\frac12 k(r_{ij} - d)^2$, with $d$ its natural
length; we take every spring to be at its natural length at the minimum,
$r_{ij} = d$, with $\vb{e} = (\vb{r}_j - \vb{r}_i)/d$ the unit vector along
it. Displace the atoms by small amounts $\vb{q}_i$ and $\vb{q}_j$. The new
separation is $(\vb{r}_j + \vb{q}_j) - (\vb{r}_i + \vb{q}_i) = d\,\vb{e} +
\vb{q}_j - \vb{q}_i$, since $\vb{r}_j - \vb{r}_i = d\,\vb{e}$, and the new
length is its length. For a vector $\vb{a}$ and a small vector $\vb{s}$ of
length $s$, expand $\lvert\vb{a} + \vb{s}\rvert^2 = (\vb{a} + \vb{s})\cdot
(\vb{a} + \vb{s})$ and set beside it the square of a guess for the length,
chosen so that the first two terms agree:
\begin{align*}
   \lvert\vb{a} + \vb{s}\rvert^2
   &= \lvert\vb{a}\rvert^2 + 2\,\vb{a}\cdot\vb{s} + s^2 ,\\
   \Bigl(\lvert\vb{a}\rvert + \frac{\vb{a}\cdot\vb{s}}{\lvert\vb{a}\rvert}
   \Bigr)^2
   &= \lvert\vb{a}\rvert^2 + 2\,\vb{a}\cdot\vb{s}
      + \frac{(\vb{a}\cdot\vb{s})^2}{\lvert\vb{a}\rvert^2} .
\end{align*}
The two squares differ by $s^2 - (\vb{a}\cdot\vb{s})^2/\lvert\vb{a}
\rvert^2$, of order $s^2$ since $\lvert\vb{a}\cdot\vb{s}\rvert = \lvert
\vb{a}\rvert s\,\lvert\cos\theta\rvert \le \lvert\vb{a}\rvert s$, with
$\theta$ the angle between them. Two positive numbers
$X$ and $Y$ differ by $(X^2 - Y^2)/(X + Y)$, and here $X + Y$ is close to
$2\lvert\vb{a}\rvert$, so the numbers themselves differ by order $s^2$
too:
\begin{equation*}
   \lvert\vb{a} + \vb{s}\rvert = \lvert\vb{a}\rvert
   + \frac{\vb{a}\cdot\vb{s}}{\lvert\vb{a}\rvert} + \order{s^2} .
\end{equation*}
With $\vb{a} = d\,\vb{e}$ and $\vb{s} = \vb{q}_j - \vb{q}_i$, the stretch
is $r_{ij} - d = \vb{e}\cdot(\vb{q}_j - \vb{q}_i) + \order{q^2}$, with $q$
the size of the displacements. Squaring, the correction of order $q^2$
appears only multiplied by $\vb{e}\cdot(\vb{q}_j - \vb{q}_i)$, of order
$q$, or by itself, so the energy, to second order in the displacements,
is
\begin{equation*}
   U = \tfrac12 k\,\bigl[\vb{e}\cdot(\vb{q}_j - \vb{q}_i)\bigr]^2 :
\end{equation*}
only the part of the relative displacement along the spring stretches it.
By the chain rule, since $\vb{e}\cdot\vb{q}_j = \sum_b
e_b(\vb{q}_j)_b$ with $(\vb{q}_j)_b$ the component $b$ of $\vb{q}_j$, $\partial U/
\partial (\vb{q}_j)_a = ke_a\,\vb{e}\cdot(\vb{q}_j - \vb{q}_i)$ and $\partial U/
\partial (\vb{q}_i)_a = -ke_a\,\vb{e}\cdot(\vb{q}_j - \vb{q}_i)$. Differentiating
once more, $\partial^2U/\partial (\vb{q}_j)_a\partial (\vb{q}_j)_b = \partial^2U/
\partial (\vb{q}_i)_a\partial (\vb{q}_i)_b = ke_ae_b$ and $\partial^2U/\partial (\vb{q}_i)_a
\partial (\vb{q}_j)_b = -ke_ae_b$. The Hessian of the spring is made of these $3\times3$ blocks,
and that of a molecule is the sum over its springs. The function
\texttt{spring\_hessian} of \texttt{mdlab} builds it, and its tests check
it against second differences of the energy.

\emph{Rigid displacements cost no energy.} Move every atom by the same
small vector $\vb{c}$: then $\vb{q}_j - \vb{q}_i = \vb{0}$ for every
spring, and $U$ does not change. Turn the molecule through a small angle
about an axis along the unit vector $\vb{n}$ through the origin: by
\cref{eq:ro-rigid}, turning at the angular speed $\omega$ for a short time
$t$ moves each atom by its velocity times $t$, $\vb{q}_i = \omega t\,\vb{n}
\times\vb{r}_i = \theta\,\vb{n}\times\vb{r}_i$, with $\theta =
\omega t$ the small angle turned. Then
\begin{equation*}
   \vb{e}\cdot(\vb{q}_j - \vb{q}_i)
   = \theta\,\vb{e}\cdot\bigl(\vb{n}\times(\vb{r}_j - \vb{r}_i)\bigr)
   = \theta d\,\vb{e}\cdot(\vb{n}\times\vb{e}) = 0 ,
\end{equation*}
since $\vb{n}\times\vb{e}$ is at right angles to $\vb{e}$. Again no spring
is stretched. The force of a spring on atom $j$ is minus the derivative of
its energy, by the chain rule $-k\,\vb{e}\,[\vb{e}\cdot(\vb{q}_j -
\vb{q}_i)]$, and that on atom $i$ is its opposite, so for both kinds of
rigid displacement every force vanishes. The force is
$-\boldsymbol{\Phi}\vb{q}$ (\cref{sec:os-modes}), so $\boldsymbol{\Phi}
\vb{q} = \vb{0}$. With $w_b = \sqrt{m_b}\,q_b$, where $a$ and $b$ now run
over all $3N$ coordinates as in \cref{sec:os-modes}, row $a$ of the
mass-weighted Hessian gives
\begin{equation*}
   \sum_b\frac{\Phi_{ab}}{\sqrt{m_am_b}}\,\sqrt{m_b}\,q_b
   = \frac{1}{\sqrt{m_a}}\sum_b\Phi_{ab}\,q_b = 0 ,
\end{equation*}
cancelling $\sqrt{m_b}$: $\vb{w}$ is an eigenvector with eigenvalue zero, a
normal mode of zero frequency (\cref{ex:ro-null} checks a case by hand).
The three drifts and three turnings are independent, none a combination
of the others, and a symmetric matrix has as many zero eigenvalues as it
has independent vectors that it turns into zero (a general result we use
without proof, like that of \cref{sec:os-modes}), so they give six zero
modes, five for a linear molecule. The other $3N - 6$ ($3N - 5$)
eigenvalues are positive only if every change of shape costs energy, as at
the minimum of a real molecule.

The same holds for any potential energy that does not change when the
whole molecule is moved or turned, as for every isolated molecule,
provided the atoms sit at a minimum. A moved or turned copy of the minimum
is then a minimum too, where every force vanishes. A small rigid
displacement $\eta\vb{u}$ lands within $\order{\eta^2}$ of
such a copy, so the force there, $-\eta\boldsymbol{\Phi}\vb{u}$ to
first order, is of order $\eta^2$; dividing by $\eta$ and
letting $\eta$ shrink leaves $\boldsymbol{\Phi}\vb{u} = \vb{0}$.

\subsection*{A spring model of water}

Model water, at the measured geometry and with the masses of
\cref{sec:ro-tensor}, by three springs: two O--H bonds of stiffness
$k_{\mathrm{OH}}$, and an H--H spring of stiffness $k_{\mathrm{HH}}$ that
stands in for the molecule's resistance to bending; without it the bend
would cost no energy and give a seventh zero eigenvalue. Its Hessian is
$9\times9$, and with the masses of the atoms \texttt{normal\_modes} gives
nine eigenvalues. Six are zero to within the rounding error of the computer's
arithmetic, which keeps about 16 significant figures, all
below \SI{1e-15}{\per\femto\second\squared} against
\SI{0.09}{\per\femto\second\squared} for the slowest vibration: the three
drifts and the three turnings. In the antisymmetric
stretch one hydrogen atom moves out as the other moves in, and the H--H
distance does not change to first order, so $k_{\mathrm{OH}}$ alone sets
its frequency. Choosing $k_{\mathrm{OH}} = \SI{48.48}{\electronvolt\per
\angstrom\squared}$ to match the measured antisymmetric stretch of
\cref{tab:os-vibrations}, and then $k_{\mathrm{HH}} =
\SI{17.66}{\electronvolt\per\angstrom\squared}$ to match the bend, makes
the other three eigenvalues 0.09028, 0.5006 and
\SI{0.7338}{\per\femto\second\squared}; with $\omega$ the square root of
each and $\tilde\nu = \omega/(2\pi c)$ from \cref{sec:os-molecules}, they
are 1595, 3756 and \SI{4548}{\per\centi\metre}, the last, the symmetric
stretch, \SI{24}{\percent} above the measured \SI{3657}{\per\centi
\metre}. The count of modes is exact; their
frequencies are only as good as the model, and three springs are a crude
model of a molecule. A single spring between two atoms likewise gives five
zero eigenvalues and one vibration.

\Needspace{8\baselineskip}
\begin{bridge}
A vibrational analysis from electronic-structure calculations builds the
Hessian, usually by finite differences of the forces, and diagonalises its
mass-weighted form (\cref{sec:os-molecules}). The six rigid modes then come
out close to zero but not at zero: the geometry is never exactly at the
minimum, and the forces left over give the turnings a small stiffness of
either sign; the finite differences add errors of their own; and a
real-space integration grid is not quite unchanged when the atoms are
moved or turned against it. The rigid motions can then be
projected out of the Hessian before it is diagonalised, the matrix
counterpart of the recipe of \cref{sec:ro-remove}.
\end{bridge}

The count returns in Chapter~12, which measures temperature by the kinetic
energy shared among the degrees of freedom in which the atoms are free to
move. A simulation that removes the drift has three fewer, and one of an
isolated molecule that also removes the turning has six fewer, five if the
molecule is linear.

\begin{keyidea}
A molecule of $N$ atoms has $3N$ degrees of freedom: three drifts, three
turnings (two if it is linear) and $3N - 6$ ($3N - 5$) internal motions.
The rigid motions stretch nothing, and appear as zero eigenvalues of the
mass-weighted Hessian.
\end{keyidea}

\notebook{05}{the nine modes of the spring model of water}

\begin{selfcheck}
How many internal degrees of freedom has ethane, C$_2$H$_6$, and how many
has acetylene, C$_2$H$_2$, whose four atoms lie on a line?
\end{selfcheck}
